\chapter{Design Patterns}
\label{ch:patterns}

\newthought{This chapter describes} the design patterns used throughout OutOf5K, with implementation examples specific to this codebase.

\section{Pattern Summary}

\begin{table}[h]
    \centering
    \begin{tabular}{@{}lll@{}}
        \toprule
        \textbf{Pattern} & \textbf{Category} & \textbf{Primary Use} \\
        \midrule
        Repository & Structural & Database access \\
        Command & Behavioral & Tauri IPC \\
        Observer & Behavioral & Events, stores \\
        Factory & Creational & Parser creation \\
        Strategy & Behavioral & Analysis algorithms \\
        Adapter & Structural & API clients \\
        Facade & Structural & Python bridge \\
        Builder & Creational & Query construction \\
        \bottomrule
    \end{tabular}
    \caption{Design patterns used in OutOf5K}
\end{table}

\section{Repository Pattern}

\marginnote{The Repository pattern abstracts data persistence, providing a collection-like interface for domain objects.}

The Repository pattern separates business logic from data access:

\begin{lstlisting}[language=Rust]
pub struct DemoRepository {
    pool: SqlitePool,
}

impl DemoRepository {
    pub async fn find_by_id(
        &self, 
        id: &str
    ) -> Result<Option<Demo>> {
        sqlx::query_as!(Demo, 
            "SELECT * FROM demos WHERE id = ?", 
            id
        )
        .fetch_optional(&self.pool)
        .await
    }
    
    pub async fn save(&self, demo: &Demo) -> Result<()> {
        // Insert logic
    }
}
\end{lstlisting}

\textbf{Benefits:}
\begin{itemize}
    \item Testability via mocking
    \item Database implementation flexibility
    \item Centralized query logic
\end{itemize}

\section{Command Pattern}

In Tauri, commands encapsulate IPC requests:

\begin{lstlisting}[language=Rust]
#[tauri::command]
pub async fn import_demos(
    state: State<'_, AppState>,
    paths: Vec<String>,
) -> Result<ImportResult, CommandError> {
    state.demo_service
        .import_demos(&paths)
        .await
}
\end{lstlisting}

Frontend invocation:

\begin{lstlisting}[language=TypeScript]
import { invoke } from '@tauri-apps/api/tauri';

const result = await invoke('import_demos', { 
    paths: ['/path/to/demo.dem'] 
});
\end{lstlisting}

\section{Observer Pattern}

\subsection{Svelte Stores}

Svelte's reactive stores implement the Observer pattern:

\begin{lstlisting}[language=TypeScript]
import { writable } from 'svelte/store';

function createDemoStore() {
    const { subscribe, update } = writable({
        demos: [],
        loading: false,
    });

    return {
        subscribe,
        addDemo(demo) {
            update(s => ({
                ...s,
                demos: [demo, ...s.demos]
            }));
        },
    };
}

export const demoStore = createDemoStore();
\end{lstlisting}

\subsection{Tauri Events}

Backend-to-frontend notifications:

\begin{lstlisting}[language=Rust]
// Backend
app.emit_all("demo:progress", payload);

// Frontend
listen('demo:progress', (event) => {
    updateProgress(event.payload);
});
\end{lstlisting}

\section{Factory Pattern}

\marginnote{The Factory pattern creates objects without specifying the exact class, useful when creation depends on runtime conditions.}

Creating the appropriate parser for demo format versions:

\begin{lstlisting}[language=Python]
class DemoParserFactory:
    @staticmethod
    def create(file_path: Path) -> DemoParser:
        version = detect_version(file_path)
        
        if version == DemoVersion.CS2:
            return CS2DemoParser(file_path)
        elif version == DemoVersion.CSGO:
            return LegacyCSGOParser(file_path)
        else:
            raise UnsupportedDemoError()
\end{lstlisting}

\section{Strategy Pattern}

Different analysis algorithms share a common interface:

\begin{lstlisting}[language=Python]
class Analyzer(ABC):
    @abstractmethod
    def analyze(self, demo: ParsedDemo) -> Result:
        pass

class StatsAnalyzer(Analyzer):
    def analyze(self, demo: ParsedDemo) -> Result:
        # Calculate statistics
        pass

class HighlightAnalyzer(Analyzer):
    def analyze(self, demo: ParsedDemo) -> Result:
        # Detect highlights
        pass

# Usage
pipeline = AnalysisPipeline([
    StatsAnalyzer(),
    HighlightAnalyzer(),
])
results = pipeline.run(parsed_demo)
\end{lstlisting}

\section{Adapter Pattern}

Adapting external APIs to internal interfaces:

\begin{lstlisting}[language=Rust]
#[async_trait]
pub trait PlayerStatsProvider {
    async fn get_stats(&self, id: &str) 
        -> Result<PlayerStats>;
}

// HLTV adapter
pub struct HLTVAdapter { /* ... */ }

#[async_trait]
impl PlayerStatsProvider for HLTVAdapter {
    async fn get_stats(&self, id: &str) 
        -> Result<PlayerStats> {
        let html = self.fetch_page(id).await?;
        let hltv_stats = parse_html(&html)?;
        
        // Adapt to internal format
        Ok(PlayerStats {
            rating: hltv_stats.rating_2_0,
            kd_ratio: hltv_stats.kd,
            // ...
        })
    }
}
\end{lstlisting}

\section{Facade Pattern}

The Python bridge provides a simplified interface to the parsing subsystem:

\begin{lstlisting}[language=Rust]
pub struct PythonBridge {
    process: Child,
    // ...
}

impl PythonBridge {
    // High-level API (Facade)
    pub async fn parse_demo(&mut self, path: &Path) 
        -> Result<ParsedDemo> {
        self.send_command("parse", json!({
            "file_path": path
        })).await
    }
    
    pub async fn get_metadata(&mut self, path: &Path) 
        -> Result<DemoMetadata> {
        self.send_command("get_metadata", json!({
            "file_path": path
        })).await
    }
}
\end{lstlisting}

\section{Builder Pattern}

Constructing complex database queries:

\begin{lstlisting}[language=Rust]
let (query, bindings) = DemoQueryBuilder::new()
    .with_status(DemoStatus::Parsed)
    .with_map("de_dust2")
    .played_after(week_ago)
    .order_by_played_at(true)
    .limit(20)
    .build();
\end{lstlisting}
